\chapter{The Five Methods}
\label{ch:five-methods}

\begin{versebox}
Who leads through the distortions' defilements ---\\
who uproots the unwholesome at its roots ---\\
having looked with direction-gazing,\\
by craving and by ignorance, through calm.
\end{versebox}

\noindent The sixteen modes are the tools. The five methods are the strategy for using them. If the modes are like a surgeon's instruments --- scalpel, clamp, suture, probe --- then the methods are the surgical plans: which instrument to use first, in what order, and why.\footnote{Pali: \textit{``Tattha katamaṃ nayasamuṭṭhānaṃ?''} ``What is the Arising of the Methods?'' The word \textit{naya} means ``method, way, approach, strategy.'' The \textit{nayasamuṭṭhāna} is the section that explains how the methods arise from the analysis of defilements. Five methods are named in the Netti's mnemonic verses: \textit{nandiyāvaṭṭa} (the Delight-Round), \textit{sīhavikkīḷita} (the Lion's Play), \textit{disālocana} (Direction-Gazing), \textit{aṅkusa} (the Goad), and \textit{tipukkhala} (the Triple-Turning). This chapter covers primarily the first three; the Goad and Triple-Turning are structural methods whose work is distributed throughout the text.}

\section*{The Two Characters}

The starting point of all five methods is a distinction so fundamental that the entire system pivots on it: every being falls into one of two character-types.

There is the view-character --- the person dominated by ignorance. And there is the craving-character --- the person dominated by desire.\footnote{Pali: \textit{``Pubbā koṭi na paññāyati avijjāya ca bhavataṇhāya ca. Tattha avijjānīvaraṇā sattā avijjāsaṃyuttā avijjāpakkhena vicaranti, te vuccanti diṭṭhicaritāti. Taṇhāsaṃyojanā sattā taṇhāsaṃyuttā taṇhāpakkhena vicaranti, te vuccanti taṇhācaritāti.''} ``No first point is discernible for ignorance and for craving for existence.'' Beings shrouded by the hindrance of ignorance wander on the side of ignorance --- they are called ``view-characters'' (\textit{diṭṭhicarita}). Beings fettered by the bond of craving wander on the side of craving --- they are called ``craving-characters'' (\textit{taṇhācarita}). The entire analytical system of the Netti rests on this binary distinction.}

When these two types go astray outside the Buddha's teaching, they go astray in opposite directions. The view-character, dominated by ignorance, falls into self-mortification --- punishing the body, believing that suffering produces purification. The craving-character, dominated by desire, falls into sensual indulgence --- chasing pleasure, believing that the world's sweetness is the point. Both are wrong. Both are looking for happiness in the wrong place. And both are wrong for a specific reason: outside the Buddha's teaching, there is no establishment of the truths, no skill in calm and insight, no attainment of the peace of liberation.\footnote{Pali: \textit{``Diṭṭhicaritā ito bahiddhā pabbajitā attakilamathānuyogamanuyuttā viharanti. Taṇhācaritā ito bahiddhā pabbajitā kāmesu kāmasukhallikānuyogamanuyuttā viharanti \ldots ito bahiddhā natthi saccavavatthānaṃ, kuto catusaccappakāsanā vā samathavipassanākosallaṃ vā upasamasukhappatti vā!''} The Netti's diagnosis: the two extremes the Buddha rejected in his first sermon (self-mortification and sensual indulgence) are not arbitrary errors. They are the \textit{predictable} errors of two specific character-types operating without the framework of the Four Noble Truths.}

Without the medicine of calm and insight, both types only make their disease worse. They increase the disease, increase the tumor, increase the dart. Beaten by disease, oppressed by the tumor, pierced by the dart, they plunge up and down through hell, the animal realm, the ghost realm, the realm of demons --- experiencing crushing blows and not finding the medicine.\footnote{Pali: \textit{``Te tadabhiññā santā rogameva vaḍḍhayanti, gaṇḍameva vaḍḍhayanti, sallameva vaḍḍhayanti \ldots rogagaṇḍasallabhesajjaṃ na vindanti.''} Three metaphors for defilement: disease (\textit{roga}), tumor (\textit{gaṇḍa}), dart (\textit{salla}). The two extremes are the disease; calm and insight are the medicine. Self-mortification and sensual indulgence are both the tumor, not the cure; the arrow, not the tweezers.}

\section*{The Twenty-fold Identity-View}

The view-character and the craving-character grasp at self in different ways. The view-character identifies each aggregate \textit{as} self: ``form is my self,'' ``feeling is my self,'' ``perception is my self.'' The craving-character identifies a self that \textit{possesses} the aggregate, or that contains it, or is contained within it: ``my self has form,'' ``form is in my self,'' ``my self is in form.''\footnote{Pali: \textit{``Diṭṭhicaritā rūpaṃ attato upagacchanti \ldots Taṇhācaritā rūpavantaṃ attānaṃ upagacchanti, attani vā rūpaṃ, rūpasmiṃ vā attānaṃ \ldots ayaṃ vuccati vīsativatthukā sakkāyadiṭṭhi.''} Twenty-fold identity-view (\textit{vīsativatthukā sakkāyadiṭṭhi}): 5 aggregates × 4 modes of self-identification = 20. The view-character's mode (``this is self'') maps to annihilationism (\textit{uccheda}); the craving-character's modes (``self has this,'' ``this is in self,'' ``self is in this'') map to eternalism (\textit{sassata}). Both are extremes; the Middle Way --- the Noble Eightfold Path --- is the counterpart of both.}

Five aggregates, four modes of identification --- twenty views in all. And the view-character's mode maps to annihilationism (because if the aggregates \textit{are} the self, then when the aggregates perish, the self perishes). The craving-character's modes map to eternalism (because if the self \textit{has} the aggregates, the self might survive their destruction). Both are extremes. Both are the cycle of existence spinning. The counterpart of both is the Middle Way.

In summary: the cycle of existence is suffering. Attachment to it is the origin. The cessation of craving is the cessation of suffering. The Noble Eightfold Path is the way to that cessation. Four truths. Every time.\footnote{Pali: \textit{``Pavatti dukkhaṃ, tadabhisaṅgo taṇhā samudayo, taṇhānirodho dukkhanirodho, ariyo aṭṭhaṅgiko maggo dukkhanirodhagāminī paṭipadā. Imāni cattāri saccāni. Dukkhaṃ pariññeyyaṃ, samudayo pahātabbo, maggo bhāvetabbo, nirodho sacchikātabbo.''} The Netti's refrain: every analysis, however elaborate, reduces to the same four truths with the same four tasks (comprehend, abandon, develop, realize).}

\section*{The Four Directions of Defilement}

And now the Netti reveals its most elaborate structural mapping --- what it calls the ``four directions.'' This is the architecture of suffering laid out like a blueprint.

Take ten categories of defilement: the four nutrients, the four distortions, the four clingings, the four yokes, the four body-ties, the four taints, the four floods, the four darts, the four consciousness-stations, and the four wrong courses. Within each category, the first two items belong to the craving-character and the last two to the view-character.\footnote{Pali: \textit{``Upacayena sabbepi kilesā catūhi vipallāsehi niddisitabbā. Te kattha daṭṭhabbā? Dasa vatthuke kilesapuñje.''} ``By way of accumulation, all defilements should be analyzed through the four distortions. Where should they be seen? In ten defilement-groups with their bases.'' The ten groups: (1) nutrients, (2) distortions, (3) clingings, (4) yokes, (5) body-ties, (6) taints, (7) floods, (8) darts, (9) consciousness-stations, (10) wrong courses. Each group has four members. The first member of each group links to the first of all the others, forming a ``direction.'' There are four such directions.}

And each category cascades into the next. Start with the first ``direction'':

Physical food, seeing beauty in the unattractive, sense-pleasure clinging, the sense-pleasure yoke, the body-tie of covetousness, the sense-pleasure taint, the sense-pleasure flood, the dart of lust, the form-based consciousness-station, the wrong course of desire. Ten terms. One meaning. Only the wording differs.\footnote{Pali: \textit{``Kabaḷīkāro āhāro `asubhe subha'nti vipallāso, kāmupādānaṃ, kāmayogo, abhijjhākāyagantho, kāmāsavo, kāmogho, rāgasallo, rūpūpagā viññāṇaṭṭhiti, chandā agatigamananti --- paṭhamā disā.''} The first direction: ten defilement-terms that are all synonyms, all describing the same syndrome from different angles. Physical food triggers the beauty-distortion, which triggers sense-pleasure clinging, which yokes one to sense-pleasure, which ties the body with covetousness, which ferments into the sense-pleasure taint, which swells into the sense-pleasure flood, which lodges as the dart of lust, which stations consciousness in form, which drives one on the wrong course of desire.}

Now the cascade makes sense: each distortion triggers the clinging that matches it, each clinging triggers the yoke, each yoke triggers the body-tie, each body-tie ferments into a taint, each taint swells into a flood, each flood lodges as a dart, each dart stations consciousness in a specific aggregate, and each consciousness-station drives a specific wrong course. It is not that there are forty separate problems. There are four problems, each with ten faces.

The second direction: contact as nutriment, seeing pleasure in suffering, existence-clinging, the existence-yoke, the body-tie of ill-will, the existence-taint, the existence-flood, the dart of hatred, the feeling-based consciousness-station, the wrong course of hatred.

The third direction: consciousness as nutriment, seeing permanence in the impermanent, view-clinging, the view-yoke, the body-tie of grasping at rules, the view-taint, the view-flood, the dart of conceit, the perception-based consciousness-station, the wrong course of fear.

The fourth direction: mental intention as nutriment, seeing self in what is not-self, self-doctrine clinging, the ignorance-yoke, the body-tie of ``this alone is true,'' the ignorance-taint, the ignorance-flood, the dart of delusion, the formation-based consciousness-station, the wrong course of delusion.\footnote{Pali: \textit{``Phasso āhāro, `dukkhe sukha'nti vipallāso \ldots dosā agatigamananti --- dutiyā disā. Viññāṇāhāro `anicce nicca'nti vipallāso \ldots bhayā agatigamananti --- tatiyā disā. Manosañcetanāhāro `anattani attā'nti vipallāso \ldots mohā agatigamananti --- catutthī disā.''} Four directions, each with ten parallel terms. The first two directions belong to the craving-character; the last two to the view-character (subdivided into ``dull'' and ``keen''). The Netti's assertion: these forty terms reduce to four syndromes, and the forty terms found scattered across the entire canon are not separate teachings but four versions of the same diagnosis.}

\section*{The Medicine for Each Direction}

And each direction has its own liberation-door as medicine. The first and second directions (the craving-character's syndromes) are abandoned through the \textit{desireless} liberation-door. The third direction (the dull view-character) is abandoned through the \textit{emptiness} door. The fourth direction (the keen view-character) is abandoned through the \textit{signless} door.\footnote{Pali: \textit{``Tattha yo ca kabaḷīkāro āhāro yo ca phasso āhāro, ime appaṇihitena vimokkhamukhena pariññaṃ gacchanti; viññāṇāhāro suññatāya; manosañcetanāhāro animittena.''} The three liberation-doors (\textit{vimokkhamukhā}) mapped onto the four directions: desireless (\textit{appaṇihita}) for directions 1--2, emptiness (\textit{suññatā}) for direction 3, signless (\textit{animitta}) for direction 4. Thus every defilement in the entire system has a specific remedy.}

\section*{The Ten Practice-Groups}

The Netti then maps ten parallel groups of wholesome qualities, each with four members, as the specific antidotes to the ten defilement-groups: the four practices counter the four nutrients. The four foundations of mindfulness counter the four distortions. The four absorptions counter the four clingings. The four abidings counter the four yokes. The four right strivings counter the four body-ties. The four marvels counter the four taints. The four resolutions counter the four floods. The four concentration-developments counter the four darts. The four happiness-sharing qualities counter the four consciousness-stations. The four boundless states counter the four wrong courses.\footnote{Pali: \textit{``Tesaṃ vikkīḷitaṃ. Cattāro āhārā tesaṃ paṭipakkho catasso paṭipadā \ldots Cattāri agatigamanāni tesaṃ paṭipakkho catasso appamāṇā.''} Ten defilement-groups × 4 = 40 defilements. Ten practice-groups × 4 = 40 antidotes. Each antidote precisely counters one defilement. The Netti calls this the ``lion's play'' (\textit{sīhavikkīḷita}) --- the ``sport'' or ``range'' of Buddhas, Paccekabuddhas, and disciples who have destroyed greed, hatred, and delusion.}

And then --- cascading development. The first practice, developed and cultivated, fulfills the first foundation of mindfulness. The first foundation fulfills the first absorption. The first absorption fulfills the first abiding. The first abiding fulfills the first right striving. The first right striving fulfills the abandonment of conceit. The abandonment of conceit fulfills the resolution of truth. The resolution of truth fulfills desire-concentration. Desire-concentration fulfills sense-restraint. And sense-restraint fulfills loving-kindness. Ten links, first to first, cascading upward.\footnote{Pali: \textit{``Paṭhamā paṭipadā bhāvitā bahulīkatā paṭhamaṃ satipaṭṭhānaṃ paripūreti \ldots indriyasaṃvaro bhāvito bahulīkato mettaṃ paripūreti.''} The cascading fulfillment runs across all four directions simultaneously, yielding a 10 × 4 grid where each column represents one direction's complete path from initial practice to final boundless state: (1) practice → mindfulness → absorption → abiding → striving → abandonment of conceit → truth-resolution → desire-concentration → sense-restraint → loving-kindness; (2) similar for hatred/compassion; (3) for conceit/sympathetic joy; (4) for delusion/equanimity.}

The same for the second, third, and fourth positions --- each cascading through its own column:

The lust-character's medicine terminates in loving-kindness. The hatred-character's medicine terminates in compassion. The conceit-character's medicine terminates in sympathetic joy. The delusion-character's medicine terminates in equanimity.\footnote{Pali: \textit{``Idaṃ rāgacaritassa puggalassa bhesajjaṃ \ldots Idaṃ dosacaritassa puggalassa bhesajjaṃ \ldots Idaṃ diṭṭhicaritassa mandassa bhesajjaṃ \ldots Idaṃ diṭṭhicaritassa udattassa bhesajjaṃ.''} The terminus of each column: column 1 (lust) → loving-kindness (\textit{mettā}); column 2 (hatred) → compassion (\textit{karuṇā}); column 3 (dull view-character) → sympathetic joy (\textit{muditā}); column 4 (keen view-character) → equanimity (\textit{upekkhā}). The four boundless states (\textit{brahmavihārā}) are not generic virtues but specific antidotes for specific defilements.}

Four characters. Four diseases. Four columns of medicine. Each column with ten cascading levels. The entire Buddhist path mapped as a grid.

\section*{The Three Types of Learner}

The Netti closes the Five Methods section by collapsing the grid one more time. Four character-types become three types of learner: the one who understands on the spot, the one who understands with elaboration, and the one who needs to be led.\footnote{Pali: \textit{``Iti kho cattāri hutvā tīṇi bhavanti --- ugghaṭitaññū vipañcitaññū neyyoti.''} Four paths of practice → three types of learner: the \textit{ugghaṭitaññū} (``understands the moment it's lifted up'' --- i.e., gets it immediately), the \textit{vipañcitaññū} (``understands when elaborated''), and the \textit{neyya} (``one who can be led'' --- needs extended guidance). The fourth type --- the \textit{padaparama} (``one for whom words are the highest'') --- is not mentioned here because such a person cannot attain liberation in this life and is therefore outside the scope of the path-method.}

The Buddha teaches the one who understands on the spot with calm alone --- briefly, gently, showing only the escape. The Buddha teaches the one who needs to be led with insight alone --- at length, intensely, showing the gratification, the danger, and the escape. The one who understands with elaboration gets the middle treatment: both calm and insight, moderate length, the danger and the escape.\footnote{Pali: \textit{``Bhagavā ugghaṭitaññussa puggalassa samathaṃ upadisati \ldots mudindriyassa bhagavā vipassanaṃ upadisati \ldots majjhindriyassa bhagavā samathavipassanaṃ upadisati.''} The Buddha's teaching method calibrated to the learner: (1) keen-faculty → calm, brief instruction, escape only, higher-wisdom training; (2) dull-faculty → insight, detailed instruction, gratification + danger + escape, higher-virtue training; (3) medium-faculty → calm and insight, moderate instruction, danger + escape, higher-mind training.}

Three become two: the craving-character and the view-character. And the defilements and purifications of each are listed in exhaustive detail --- the craving-character's corrupt side runs from craving itself through heedlessness, sense-restraint failure, covetousness, ill-will, the hindrances, the fetters, anger, grudge-bearing, contempt, insolence, envy, stinginess, deceit, fraud, eternalism, and annihilationism. The purification side runs from calm and insight through conscience and prudence, mindfulness and clear comprehension, wise attention, energy, being easy to admonish, knowledge of the teaching, knowledge of inference, knowledge of destruction, knowledge of non-arising, faith, heedfulness, hearing the true teaching, restraint, non-covetousness, non-ill-will, liberation of the heart through fading of lust, liberation through wisdom by fading of ignorance, full penetration, fewness of wishes, contentment, non-anger, non-grudge-bearing, non-contempt, non-insolence, the abandoning of envy, the abandoning of stinginess, true knowledge, liberation, the conditioned liberation, the unconditioned liberation, the nirvana-element with remainder, and the nirvana-element without remainder.\footnote{Pali: \textit{``Tesaṃ dvinnaṃ puggalānaṃ ayaṃ saṃkileso --- taṇhā ca avijjā ca ahirikañca anottappañca \ldots ucchedadiṭṭhi ca. Tesaṃ dvinnaṃ puggalānaṃ idaṃ vodānaṃ --- samatho ca vipassanā ca hirī ca ottappañca \ldots anupādisesā ca nibbānadhātu.''} The comprehensive lists of defilement and purification for the two character-types. The defilement list has over 30 terms; the purification list has over 35. Every term in both lists has been encountered somewhere in the sixteen modes --- the Five Methods section is demonstrating that the system is complete: every defilement named in the canon can be located in the four-direction framework, and every antidote is accounted for.}

Two become one: calm and insight, working together, are the single medicine for the single disease of continued existence.

This is what the Venerable Mahākaccāyana meant: ``Who leads through the distortions' defilements --- who uproots the unwholesome at its roots.'' That is the domain of the Five Methods.\footnote{Pali: \textit{``Ayaṃ vuccati tipukkhalassa ca nayassa aṅkusassa ca nayassa bhūmīti. Tenāha `yo akusale samūlehi netī'ti `oloketvā disalocanenā'ti ca.''} The closing formula identifies the domain of two of the five methods: the Triple-Turning (\textit{tipukkhala}) and the Goad (\textit{aṅkusa}). The five methods are not five separate procedures but five aspects of a single strategic approach: (1) the Delight-Round analyzes the cycle of craving; (2) the Lion's Play maps the symmetry between defilement and purification; (3) Direction-Gazing surveys the four directions; (4) the Goad compresses the analysis to its essentials; (5) the Triple-Turning applies the three liberations.}
